\documentclass[11pt]{article}
\usepackage{amsmath,amssymb,booktabs,geometry,xcolor,hyperref}
\geometry{a4paper,top=24mm,bottom=24mm,left=25mm,right=25mm,headheight=14pt}
\definecolor{linkink}{RGB}{19,73,85}
\hypersetup{colorlinks=true,urlcolor=linkink,citecolor=linkink,linkcolor=linkink,
 pdftitle={A rotation-class completion criterion for seven-symbol superpermutations},
 pdfauthor={Jaehun Jung},pdfsubject={Completion thresholds for seven-symbol superpermutations}}
\urlstyle{same}
\newtheorem{theorem}{Theorem}[section]
\newtheorem{lemma}[theorem]{Lemma}
\newtheorem{proposition}[theorem]{Proposition}
\newtheorem{corollary}[theorem]{Corollary}
\newenvironment{proof}[1][Proof]{\par\smallskip\noindent\textit{#1.}\enspace}{\unskip\nobreak\hfill\ensuremath{\square}\par\smallskip}
\setlength{\emergencystretch}{2em}
\pagestyle{myheadings}
\markright{\small J. Jung\hfill Rotation-class completion on seven symbols\hspace{1.2em}}
\title{A rotation-class completion criterion for seven-symbol superpermutations}
\author{Jaehun Jung\\[4pt]{\small Independent Researcher, Republic of Korea}}
\date{24 November 2026}
\begin{document}
\maketitle
\begin{abstract}
We prove that $M_k+7(720-k)=L(7)$ for every integer $91\le k\le720$, where $L(7)$ is the minimum length of a seven-symbol superpermutation and $M_k$ is the minimum length of a word that visits all 720 rotation classes and completes at least $k$. A six-letter factor with distinct symbols determines the class it visits. Completion requires all seven cyclic shifts to occur. At $k=91$, a superpermutation of length at most 5,904 exists exactly when a word of length 1,501 visits every class and completes 91. Our examples visit all 720 classes, completing 90 at length 1,501 and 91 at length 1,502. Halfar's computer-assisted lower bound and Teodorescu's word W-0021 give $5901\le L(7)\le5905$. We prove that 5,905 is optimal for words retaining a prescribed 463-edge residual route and at least 63 six-edge wheels, and for words with the exceptional support of any of 98 record classes in Halfar's catalogue. Whether 5,905 is optimal over all words remains open.
\end{abstract}
\noindent\textit{Keywords.} Superpermutations, rotation classes, Euler trails, shortest common superstrings, computer-assisted proof.

\section{Introduction}
A superpermutation on an alphabet $A$ contains each permutation of $A$ as a contiguous factor. Write $L(n)$ for its minimum length on $n$ symbols. Ashlock and Tillotson~\cite{ashlock} obtained the lower bound $n!+(n-1)!+n-2$ by counting the letters needed to move between cyclic classes. Engen and Vatter~\cite[\S2]{engenvatter} explain this count and give the anonymous lower-bound argument that also charges changes of trajectory. It gives $n!+(n-1)!+(n-2)!+n-3$. Both counts apply to a word that completes every class. Houston~\cite{houston} used a different formulation, a weighted overlap graph on all permutations, to construct a six-symbol word of length 872.

For $n\ge5$, Halfar's defect formulation~\cite{halfardefect} relates the existence of a full superpermutation at a prescribed length to that of a configuration of the corresponding defect. For seven symbols, his computer-assisted lower bound~\cite{halfar} and Teodorescu's word W-0021~\cite{teodorescu} give
\begin{equation}
       5901\leq L(7)\leq 5905. \label{eq:sevenbounds}
\end{equation}
The 5,901 endpoint uses external finite searches in addition to a Lean reduction. The proof of the threshold identity below uses W-0021 and is independent of those external searches.

Here a partial word must reach all 720 rotation classes, but the number it must complete can vary. The distinction can already be seen at the six-letter context $012345$. Its missing symbol is $6$, so the context identifies the class of $0123456$ even when the next letter is another $5$. Shorthand universal cycles similarly identify a permutation by appending its missing symbol to an $(n-1)$-letter factor~\cite{shorthand}. They require every such factor to occur in a circular word. For a visit in our linear word, any one of the seven contexts belonging to a class suffices. To complete that class, all seven rotations must occur as actual seven-letter factors.

Let $M_k$ be the minimum length of a word that visits every class and completes at least $k$. Inserting $6012345$ at the end of the context $012345$ completes its class and returns to the same context. The corresponding insertion at each incomplete class gives $L(7)\le M_k+7(720-k)$. The substantive direction is to recover a partial word from a shortest full word with the same saving. Theorem~\ref{thm:sevencurve} proves this for $91\le k\le720$ by bounding the number of cycles needed to join the residual components. Thus the equality concerns the minimum over all partial words. Section~\ref{sec:decision} applies it at $k=91$ to obtain the decision at length 1,501.

\section{Rotation classes and the residual graph}
\label{sec:graph}


Fix the alphabet $A=\{0,\ldots,6\}$. The seven cyclic shifts of a
permutation form a \emph{rotation class}. There are $7!/7=720$ classes.
A six-letter factor $u$ with distinct symbols has a unique missing symbol
$x$ and hence a label, the rotation class of $ux$. A word \emph{visits}
this class when $u$ occurs, including at the right end of the word.
It \emph{completes} the class when each of its seven permutations occurs
as a contiguous factor. Let $V(v)$ and $C(v)$ count the classes visited
and completed by $v$, and put
\[
 M_k=\min\{\,|v|: V(v)=720,\ C(v)\ge k\,\}.
\]
The minimum exists since W-0021 qualifies for every $0\le k\le720$.

A seven-letter factor is an edge from its six-letter prefix to its
six-letter suffix. A word of length $\ell\ge6$ spells a walk of
$\ell-6$ edges. Let $D$ contain each of the $5{,}040$ permutation edges
once. These edges split into the $720$ directed seven-cycles of the
rotation classes. Deleting $D$ from the edge multiset of a superpermutation
$w$ leaves a multiset $X$ of size $|w|-5046$. For initial and terminal
contexts $s,t$, its degrees satisfy
\[
  \operatorname{out}_X(u)-\operatorname{in}_X(u)
  =\mathbf1_{u=s}-\mathbf1_{u=t}.
\]
Call $X$ \emph{feasible} for $s,t$ when this identity holds and every
edge in $D+X$ has its source reachable from $s$ in $D+X$. These conditions
permit an Euler walk of $D+X$ from $s$ to $t$ and hence a word with precisely
those edge multiplicities. In particular $|w|\le5904$ gives $|X|\le858$.

We use two different cycles in this graph. At a visited class with context $u$ and missing symbol $x$, appending $xu$ completes its seven-cycle and returns the walk to $u$. A \emph{wheel} is the six-edge
circuit obtained by appending the six distinct symbols of $u$ in their
original order. Its six rotated contexts label six different classes.
The seven-cycle is the operation counted by $7(720-k)$ below.

\section{Completion thresholds}
\label{sec:thresholds}

\begin{lemma}[Residual component bound]\label{lem:sevencap}
If $X$ is feasible for $s,t$ and $|X|\le858$, at most $91$ weak
components of $X$ outside the component of $s$ contain an edge.
\end{lemma}

\begin{proof}
Let $b$ be the number of noninitial components with an edge and $B$ the
number of these with at least seven edges. For each component let $c$ be
the number of distinct rotation labels it meets, and put
$r=\sum\max(c-1,0)$, with the sum over all components, including the one containing $s$. Connecting the 720
labels requires $r\ge719$, since a component meeting $c$ labels can
supply at most $\max(c-1,0)$ independent joins. Write $e$ for a component's edge count and $h=\max(c-1,0)$. The initial route has $45h\le38e+59$. Each noninitial component containing an edge has at least six edges and satisfies
\[
  45h+3+22\,\mathbf1_{e\ge7}\le38e.
\]
These are the open-route and closed-circuit estimates combined in \texttt{VirtualComponentPrice}. Summing them and using $|X|\le858$ gives
\[
  3b+22B\le38\cdot858+59-45\cdot719=308.
\]
Suppose $b\ge92$.
Then $B\le1$. Every noninitial component with fewer than seven edges is a
six-edge wheel. Indeed, the first edge entering such a component from outside in $D+X$ must belong to $D$, so its endpoint has six distinct symbols. The component is balanced and has an Euler circuit based there. Returning to that context takes at least six steps, and at six steps all six appended symbols are prescribed.

\smallskip\noindent\emph{At least 92 wheels.} Suppose there are at least 92 such components. This includes every case with $b\ge93$, as well as $b=92$, $B=0$. Set aside any 92 of these wheels. Their complement is a single route by \texttt{pure\_92} in \texttt{BaseWheelFrontier92}. The route has at most $858-6\cdot92=306$ edges. To connect all labels it needs at least $720-5\cdot92=260$ labels of its own, since each attached wheel adds at most five. For a route of $e$ edges with $c$ labels the open-route estimate gives $45c\le38(e+1)+66$. These bounds force $e=306$ and $c=260$. Starting from these labels,
one can adjoin a wheel as soon as it meets a label already reached.
Feasibility requires this procedure to reach all 720 labels. Since
$260+5\cdot92=720$, every added wheel must supply five new labels.

Here the finite test concerns the actual label sets. A normalized
enumeration of the admissible route labels is certified in \texttt{NormalizedClosed306} and \texttt{FusedPacked92}. The lookup table has 47,603 entries. Forced-wheel certificates exclude 46,958 entries by the following rule. For a specified
missing label, every wheel containing it except possibly one meets the
reached set in at least two labels. A wheel with two such contacts would add at most four new labels. A cover must therefore use the sole remaining wheel, if one exists.
The other 645 entries reduce by symmetry to 322 representatives with
checked LRAT refutations. Thus no listed set can reach all labels with
at most 92 wheels (\texttt{Cap92Lookup}).

\smallskip\noindent\emph{One long circuit and 91 wheels.} The remaining case has $b=92$ and $B=1$. Let $u,v$ be the edge counts of the initial route and the long circuit, and $i,j$ their respective label counts.
The 91 wheels use $546$ edges and supply at most $455$ joins. Thus $u+v\le312$ and $(i-1)+(j-1)\ge719-455=264$. The open-route and closed-circuit label bounds at these lengths leave the six profiles below (\texttt{VirtualLongProfiles}). For $v\ge105$ the circuit bound $71j\le60v$ excludes the first four rows by the indicated excess.
\begin{center}
\begin{tabular}{@{}rr@{\qquad}rr@{\qquad}l@{}}\toprule
\multicolumn{2}{c}{Initial route} & \multicolumn{2}{c}{Long circuit} & \\
Edges $u$ & Labels $i$ & Edges $v$ & Labels $j$ & Test \\ \midrule
104 & 90  & 208 & 176 & $71j-60v=16$ \\
130 & 112 & 182 & 154 & $71j-60v=14$ \\
156 & 134 & 156 & 132 & $71j-60v=12$ \\
182 & 156 & 130 & 110 & $71j-60v=10$ \\
284 & 242 &  28 &  24 & Period constraint \\
277 & 236 &  35 &  30 & Connected-cover test \\ \bottomrule
\end{tabular}
\end{center}
At length 28, mark one occurrence per distinct label. Saturation at 24 labels would force six marked positions
in every seven-step interval. The resulting periodic pattern forces at
least 30 distinct labels (\texttt{ShortWheelPeriod}).

In the remaining profile the 236 initial labels form one of six sets $R$
after normalization. The circuit labels form a set $F$ of 30 classes
obtained by fixing the cyclic order of five of the seven symbols. Some $p\in F$ must be
reachable from $R$ using ordinary wheels before the circuit is used.
Thus those wheels alone must span from $R\cup(F\setminus\{p\})$
(\texttt{LargeRowParent}). This enlarged starting set has at most $236+29=265$ labels and needs all 265 since $265+5\cdot91=720$. There are $\binom{7}{5}(5-1)!=504$ possible sets $F$. Among the $6\cdot504$ pairs $(R,F)$, 64 intersect in exactly one label, 218 are disjoint, and the other 2,742 share at least two labels. For a pair with one common label, only that choice of $p$ gives 265 starting labels. For disjoint sets, each of the 30 choices attains this size. Thus $64+218\cdot30=6{,}604$ triples attain the necessary size. Forced-wheel certificates exclude 6,024. The other 580 reduce to 290
representatives with checked LRAT refutations (\texttt{Mixed277}).
Hence no such starting set permits a connected cover by 91 or fewer wheels. This exhausts
the possibilities for $b\ge92$.
\end{proof}

\begin{theorem}[Minimum length at a completion threshold]\label{thm:sevencurve}
For $91\le k\le720$,
\begin{equation}
  M_k+7(720-k)=L(7).\label{eq:sevencurve}
\end{equation}
\end{theorem}

\begin{proof}
Given $v$ in the defining set for $M_k$, insert a seven-cycle at an
occurrence of a labelled context for each incomplete class. Each insertion
returns to the original context and preserves the previous walk, so the
resulting word is a superpermutation of length at most
$|v|+7(720-k)$.

For the reverse inequality take a shortest superpermutation $w$. First
suppose $|w|\le5904$. Feasibility guarantees that each class has a
labelled context among $s$ and the endpoints of the edges of $X$. Indeed,
in $D+X$ its mandatory cycle is reachable, and only a residual edge or
the initial context can first enter it. Start at the component of $s$ in
$X$. As long as a residual component remains disconnected from it, choose
a mandatory seven-cycle meeting both the reached part and a fresh
residual component. Such a cycle exists because $D+X$ is reachable and
different mandatory cycles have disjoint context sets. Add that cycle.
Each addition reaches another residual component with an edge, so
Lemma~\ref{lem:sevencap} requires at most 91 selected cycles. The initial
component of the augmented graph now contains an anchor for every class,
as well as every edge of each selected cycle. Its degree balance permits an Euler walk. Suppose $j$ cycles were selected. Spelling the walk uses at most $6+|X|+7j$ letters and completes at least those $j$ classes. Insert at most $k-j$ further seven-cycles to reach $k$ completions. Its length is at most
\[
  6+|X|+7k=|w|-7(720-k).
\]

If $|w|>5904$, equation~\eqref{eq:sevenbounds} gives $|w|=5905$.
Deleting 653 designated nonoverlapping seven-edge laps from W-0021
gives an explicit word of length 1,334 visiting all classes and
completing 67. The theorem \texttt{ComponentCurve.reference653} verifies these
three counts. Inserting $k-67$ cycles gives the required upper bound
$1334+7(k-67)=5905-7(720-k)$. This proves~\eqref{eq:sevencurve}
in both cases.
\end{proof}

\begin{corollary}[Cost of one additional completion]\label{cor:marginal}
For $91\le k<720$, one has $M_{k+1}-M_k=7$.
\end{corollary}
\begin{proof}
Subtract the identity of Theorem~\ref{thm:sevencurve} at $k$ from the identity at $k+1$ and rearrange.
\end{proof}
The two minima allow different words. The seven-letter difference therefore holds even when the entire word is chosen anew for the higher threshold.

\section{A one-letter boundary for 91 completions}
\label{sec:decision}
\begin{corollary}\label{cor:decision}
A seven-symbol superpermutation of length at most 5,904 exists if and only if a word of length exactly 1,501 visits every rotation class and completes 91. Among words of that length with all visits, the maximum number of completions is either 90 or 91.
\end{corollary}
\begin{proof}
Putting $k=91$ in~\eqref{eq:sevencurve} gives $L(7)=M_{91}+4403$.
In particular, $1498\le M_{91}\le1502$ by~\eqref{eq:sevenbounds}.
At $k=108$ the same theorem gives $M_{108}=M_{91}+119$.
For the following witnesses, read W-0021 from left to right. At each noninitial position $i$ where the next thirteen letters have the form $uxu$ with $u$ a six-letter context of distinct symbols and $x$ its missing symbol, mark the first seven positions. There are 653 such starts and the marked blocks are disjoint. Delete the first 629 blocks.
\texttt{Core1501.controls} checks both words against the definitions.
\begin{center}
\begin{tabular}{@{}lrrr@{}}\toprule
Construction & Length & Visited & Completed \\ \midrule
Delete 629 laps & 1,502 & 720 & 91 \\
Delete the first letter of that word & 1,501 & 720 & 90 \\
\bottomrule\end{tabular}
\end{center}
A word of length at most 1,501 visiting all classes cannot complete 92. Inserting at most $720-92$ seven-cycles would give a superpermutation of
length at most $1501+7\cdot628=5897$, already below the formally
verified lower bound 5,899 \cite{halfar}. Repeating a word's first
letter at its beginning adds one six-context and one seven-factor
containing a repeated letter, and preserves both counts. Thus the maximum
number of completed classes at length 1,501 among words visiting all classes
is either 90 or 91. In the former case $L(7)=5905$. In the latter case a
word of length at most 5,904 exists. It follows that $L(7)=5905$ if and only if every word of length at most 1,620 that visits all classes completes at most 107.

\end{proof}

\section{Two families with minimum length 5,905}
\subsection{A prescribed residual route}
The route \texttt{ResidualControls.mainWord} in W-0021 uses 463 edges and visits 390 classes. Fix its edge multiset. Among
superpermutations whose residual edges contain this multiset together
with at least 63 whole wheels, the exact minimum is 5,905. Containment
counts multiplicities and places no requirement on the order in which
the walk reads the prescribed edges. With 63 selected wheels the least
number of additional residual edges is 18, and with 64 it is 12. Both budgets give the same length, since
$6+5040+463+6\cdot63+18=5905$ and $6\cdot64+12=6\cdot63+18$.
For the 63-wheel lower bound the budget at length 5,904 leaves 17
additional edges. The degree and connectivity constraints force a circuit
among them. A loop can be removed. Another wheel reduces the question to
the 64-wheel bound. For a nonwheel circuit with at least two edges, a count
of rotation labels forces its length to be at most 12, leaving at least
five edges that must extend the ends of the fixed route. The extension
loses a further label, contradicting coverage. These alternatives are
formalized in \texttt{Reference63.no\_seventeen}. The matching witness
and the case of more wheels give \texttt{Reference63.restricted\_minimum}.

\subsection{Exceptional supports from 98 recorded words}
Let $P$ send each permutation to its successor in a cyclic ordering of the 5,040 permutations, and let $R$ be left rotation. In the cycle decomposition of $R^{-1}P$, remove the fixed points and the ordinary cycles that hold the first symbol fixed while rotating the other six. The remaining permutations form the exceptional support. For each of the 98 archived record classes in Halfar's catalogue~\cite{ties}, keeping this support fixed still forces length at least 5,905, even when its elements are reordered or split into several cycles, the ordinary cycles are changed, and the position where the order is opened into a word is chosen freely. The record itself attains the bound in each case.

The finite certificates assign a departure price and an arrival price to each exceptional permutation, with a further unit charged on exiting a designated group. Each pair of consecutive elements satisfies its price inequality. A group with no exit contains an entire cycle, whose charge is already present in the length formula. Summation gives a single lower bound for all orders on that support. The supplement includes the 98 price tables, the Lean check \texttt{ArchiveOrders.all\_cycle\_partitions}, and a script that matches the literal supports to the source words. An arbitrary short word may have a different exceptional support.

\section{Excluding an entire prefix}
The two operations counted by a partial word obey a useful local restriction. For a word $u$ of length at least six, write $P(u)$ for the number of distinct seven-symbol permutations occurring in $u$. The following test uses no advance knowledge of which classes will be completed.
\begin{proposition}\label{prop:prefix}
If $ur$ visits every rotation class and completes at least 91, then
\[
                    V(u)+P(u)+|r|\ge 1{,}357.
\]
\end{proposition}
\begin{proof}
Fix any 91 classes completed by $ur$. They require $7\cdot91=637$ distinct permutation factors. Give one unit of credit for the first visit to each of the 720 classes and one for the first occurrence of each of these 637 permutations. Extending $u$ by one letter creates at most one new six-letter context and one new seven-letter factor. If the new factor is a permutation, its prefix and suffix contexts label the same rotation class, so the suffix cannot be its first visit. Each new letter therefore earns at most one credit. The prefix has earned at most $V(u)+P(u)$ and the completed word has earned $720+637$.
\end{proof}
For example, the first 1,469 letters of the length-1,501 control in Section~\ref{sec:decision} satisfy $V(u)=716$ and $P(u)=608$. Completing 91 classes while visiting all 720 requires at least $1357-716-608=33$ further letters. This excludes every one of the $7^{32}$ length-1,501 extensions of this prefix. The prefix and the bound are checked in \texttt{PrefixCharge}. 

\section{A certified encoding of the remaining decision}\label{sec:sat}
The fixed-length decision admits a Boolean encoding with $1{,}462{,}260$
variables and $5{,}679{,}593$ clauses. For each six-context the circuit
records its rotation label when its symbols are distinct. A following
permutation factor contributes its terminal symbol as a mark. For a given
class the seven possible marks distinguish its seven permutations. The
final six-context contributes a visit. The circuit groups records by label,
combines their marks, and tests both thresholds. The Lean theorem
\texttt{counterexample\_all\_rows} in \texttt{CompactCounterexample.lean} constructs a satisfying assignment
for every qualifying word. This implication suffices to deduce $L(7)=5905$
from a checked refutation of the whole formula. Two separately verified
symmetries allow first-occurrence ordering of the alphabet and simultaneous
fixing of 1,023 routing switches, respectively. No checked refutation of the full instance is available, so the 1,501-letter decision remains open. The original instance has SHA-256 digest
\begin{center}\small\ttfamily
dd6e2ee06c613845e722b0df3bd57d01ad188a7d13ca9b42a60c488fdcd015bd
\end{center}



\section{Conclusion and reproducibility}
The equality $M_{91}+4403=L(7)$ specifies what either outcome would establish. A length-1,501 witness with 91 completions would yield a superpermutation of length at most 5,904 by completing the remaining classes. A universal upper bound of 90 completions at that length would make W-0021 optimal. The verified examples attain 90 at length 1,501 and 91 one letter later.

The supplementary archive includes the Lean source for the component bound and the completion theorem, the 98 finite price certificates, instructions for rebuilding the Boolean instance, and a direct Python check of the witness words. The main equality is formalized in \texttt{StoppingCriterion.minimum\_length}. The bound of 91 completions uses the formally verified 5,899 lower bound for $L(7)$. Finite checks in the development use \texttt{native\_decide}, and certain finite subinstances have separately checked LRAT refutations. The verification guide identifies the Lean toolchains and the pinned upstream libraries.

\paragraph{Author contributions and use of AI.}
Jaehun Jung set the research questions, directed the sequence of investigations, and specified the standard of verification. OpenAI Codex (GPT-6) was used to develop proof arguments, write Lean and verification code, conduct computational experiments, and draft the manuscript. The constructions, external bounds, and formal libraries used here are credited in the references and the supplement.

\begin{thebibliography}{9}
\bibitem{ashlock} D. A. Ashlock and J. Tillotson, Construction of small superpermutations and minimal injective superstrings, \emph{Congressus Numerantium} \textbf{93} (1993), 91--98.
\bibitem{engenvatter} M. Engen and V. Vatter, Containing all permutations, \emph{American Mathematical Monthly} \textbf{128} (2021), 4--24. \href{https://doi.org/10.1080/00029890.2021.1835384}{Published article}, \href{https://arxiv.org/abs/1810.08252}{preprint}.
\bibitem{shorthand} A. E. Holroyd, F. Ruskey and A. Williams, Shorthand universal cycles for permutations, \emph{Algorithmica} \textbf{64} (2012), 215--245. \href{https://doi.org/10.1007/s00453-011-9544-z}{Published article}.
\bibitem{halfardefect} J. Halfar, $s(7)\ge5{,}899$, proved in Lean, Superpermutators research announcement (October 2026). \href{https://groups.google.com/g/superpermutators/c/dmBTmhy0a1A}{Statement of the defect characterization and proof sources}.
\bibitem{houston} R. Houston, \emph{Tackling the Minimal Superpermutation Problem}, arXiv:1408.5108 (2014). \href{https://arxiv.org/abs/1408.5108}{Preprint}.
\bibitem{halfar} J. Halfar, \emph{A superpermutation on seven symbols has at least 5,901 letters}, proof and certificate inventory, revision \texttt{8e3c02b} (2026). \href{https://github.com/jhalfar/superpermutations/tree/8e3c02b5234b49a1422751f86f9e472706082374/proofs/n7_5901}{Proof and computation} and \href{https://github.com/jhalfar/superpermutation-almanac-lean/tree/a09f9ef4ebb74ae1297f17c5fad38296695f2959}{Lean library}.
\bibitem{ties} J. Halfar, \emph{Seven-symbol record classes}, catalogue and source inventory (2026). \href{https://github.com/jhalfar/superpermutations/tree/8e3c02b5234b49a1422751f86f9e472706082374/ties}{Catalogue}.
\bibitem{teodorescu} A. Teodorescu, W-0021, \emph{Superpermutation Almanac}, 17 August 2026. \href{https://superpermutations.com/w/W-0021}{Word and coverage statistics}.
\end{thebibliography}
\end{document}
